\documentclass[10pt]{article}
\usepackage[T1]{fontenc}
\usepackage[utf8]{inputenc}
\usepackage[french]{babel}
\usepackage{amsmath,amssymb,mathrsfs,mathpazo}
\usepackage{geometry}
\geometry{a4paper,margin=2cm}
\usepackage{graphicx,array,tabularx,tabularray,colortbl,multirow,multicol,booktabs,enumitem,pifont,ulem,fancybox,diagbox,cancel}
\usepackage{pstricks,pstricks-add,pst-plot,pst-tree,pst-node,pst-text,pst-func,pst-math,pst-bspline,pst-3dplot,pst-eucl,pst-solides3d}
\let\pscalculate\undefined
\usepackage{pst-bezier}
\usepackage{tikz,pgfplots,tkz-tab}
\usetikzlibrary{arrows,arrows.meta,patterns,calc,positioning,trees,decorations.markings,shapes.geometric,shapes.misc,shapes.arrows,intersections,decorations.pathmorphing,decorations.pathreplacing,backgrounds,shadows,babel}
\pgfplotsset{compat=1.18}
\usepgfplotslibrary{fillbetween}
\usepackage[np]{numprint}
\usepackage{listings,makecell,pas-tableur,twoopt}
\let\barre\undefined
\usepackage[pstricks]{tablvar}
\definecolor{darkgreen}{RGB}{0,100,0}
\providecommand{\red}{\color{red}}
\providecommand{\blue}{\color{blue}}
\providecommand{\green}{\color{green}}
\providecommand{\black}{\color{black}}
\providecommand{\white}{\color{white}}
\providecommand{\gray}{\color{gray}}
\providecommand{\cyan}{\color{cyan}}
\providecommand{\magenta}{\color{magenta}}
\providecommand{\yellow}{\color{yellow}}
\providecommand{\violet}{\color{violet}}
\providecommand{\orange}{\color{orange}}
\providecommand{\euro}{\texteuro}
\providecommand{\EUR}{\texteuro}
\providecommand{\ieme}{\textsuperscript{e}}
\providecommand{\degres}{\ensuremath{^\circ}}
\providecommand{\degree}{\ensuremath{^\circ}}
\providecommand{\vv}[1]{\overrightarrow{#1}}
\providecommand{\Coord}[1]{#1}
\providecommand{\decosix}{}
\providecommand{\starredbullet}{\ensuremath{\bullet}}
\pagestyle{empty}
\setlength{\parindent}{0pt}
\usepackage{adjustbox}\DeclareUnicodeCharacter{2236}{:}\DeclareUnicodeCharacter{2190}{\ensuremath{\leftarrow}}
\providecommand{\R}{\mathbb{R}}
\providecommand{\N}{\mathbb{N}}
\providecommand{\Z}{\mathbb{Z}}
\providecommand{\Q}{\mathbb{Q}}
\providecommand{\D}{\mathbb{D}}
\providecommand{\C}{\mathbb{C}}
\providecommand{\e}{{\mathrm{e}}}
\providecommand{\vect}[1]{\overrightarrow{#1}}
\providecommand{\vectt}[1]{\overrightarrow{\mathrm{#1}}}
\providecommand{\vvt}[1]{\overrightarrow{\mathrm{#1}}}
\providecommand{\np}[2][]{{#2~#1}}
\providecommand{\vv}[1]{{\overrightarrow{#1}}}
\renewcommand{\d}{{\mathrm{d}}}
\providecommand{\degres}{{^\circ}}
\providecommand{\degre}{{^\circ}}
\providecommand{\blue}{\color{blue}}
\providecommand{\red}{\color{red}}
\providecommand{\enskip}{\quad}
\providecommand{\llbracket}{\lbrack\!\lbrack}
\providecommand{\rrbracket}{\rbrack\!\rbrack}
\providecommand{\euro}{{\text{€}}}
\providecommand{\no}{{\text{n°}}}
\providecommand{\interval}[2]{{[#1;#2]}}
\providecommand{\intervFO}[2]{{[#1;#2[}}
\providecommand{\dotfill}{\ldots}

\providecommand{\ds}{\displaystyle}
\providecommand{\textding}[1]{#1}
\providecommand{\Oij}{\(\mathrm{O};\vec{\imath},\vec{\jmath}\)}
\providecommand{\Oijk}{\(\mathrm{O};\vec{\imath},\vec{\jmath},\vec{k}\)}
\providecommand{\Ouv}{\(\mathrm{O};\vec{u},\vec{v}\)}
\providecommand{\intervFF}[2]{\ensuremath{[#1;#2]}}
\renewcommand{\euro}{\ifmmode\text{\texteuro}\else\texteuro\fi}
\DeclareUnicodeCharacter{00B0}{\ifmmode^{\circ}\else\textdegree\fi}
\DeclareUnicodeCharacter{20AC}{\ifmmode\text{\texteuro}\else\texteuro\fi}
% Énoncés : Éducation nationale. Transcription et corrigés : APMEP, Denis Vergès et contributeurs. Découpage et classement : CoopMaths, Sébastien Lozano, Éric Elter et contributeurs. Adaptation de la présentation : Maths Explorer.
% CC BY-SA : https://coopmaths.fr/www/about/
% https://coopmaths.fr/alea/static/bac/2021/tex/bac_2021_06_metropole_7_cor.tex
\begin{document}
\noindent\textit{APMEP / CoopMaths. CC BY-SA. Présentation adaptée par Maths Explorer.}\par\medskip
\textbf{Commun à tous les candidats}

\medskip

Dans cet exercice, les résultats des probabilités demandées seront, si nécessaire, arrondis au millième.


























\medskip

\begin{enumerate}
\item 
	\begin{enumerate}
		\item On traduit la situation par un arbre pondéré.
		
\begin{center}
\pstree[treemode=R,nodesepB=3pt,levelsep=2.75cm]{\TR{}}
{\pstree{\TR{$M$~~}\taput{0,4}}
	{\TR{$T$}\taput{0,9}
	\TR{$\overline{T}$}\tbput{0,1}
	}
\pstree{\TR{$\overline{M}$~~}\tbput{0,6}}
	{\TR{$T$}\taput{0,15}
	\TR{$\overline{T}$}\tbput{0,85}
	}	 
}		
\end{center}

\medskip
	
		\item 
Il faut trouver $P(M \cap T) = P(M) \times P_M(T) = 0,4 \times 0,9 = 0,36$.
		\item 
On a de même $P\left(\overline{M} \cap T\right) = P\left(\overline{M}\right) \times P_{\overline{M}}(T) = 0,6 \times  0,15 = 0,09$.

D'après la loi des probabilités totales :

$P(T) = P(M \cap T) + P\left(\overline{M} \cap T\right) = 0,36 + 0,09 = 0,45$.
		\item 
		Il faut trouver $P_T(M) = \dfrac{P(M \cap T)}{P(T)} = \dfrac{0,36}{0,45} = \dfrac{36}{45} = \dfrac{9 \times 4}{9 \times 5} = \dfrac{4}{5} = \dfrac{8}{10} = 0,8$.
	\end{enumerate}
\item 


	\begin{enumerate}
		\item 
On suppose que le nombre de chats est assez important pour que l'on puisse assimiler le choix des 20 chats à un tirage avec remise.
		
La variable $X$ suit donc une loi binomiale de paramètres $n = 20$ et de probabilité $p = 0,45$ trouvé à la question \textbf{1. c.}.
		\item 
On a $p\left(X = 5\right) = \binom{20}{5} \times 0,45^5 \times (1 - 0,45)^{20-5} = \np{15504}\times 0,45^5  \times 0,55^{15} \approx \np{0,0365}$ soit environ 0,037.

		\item 
La calculatrice donne $P(X < 9) \approx 0,414$.
		\item 
On sait que l'espérance $E = n \times p = 20 \times 0,45 = 9$.
		
Cela signifie que sur un grand nombre d'échantillons il y aura en moyenne 9 chats positifs par échantillon de 20.
	\end{enumerate}		
\item 
	\begin{enumerate}
		\item 
On a encore une loi binomiale de paramètres $n$ et de probabilité d'être positif de $0,45$.

On a $P(X = 0) = \binom{n}{0}\times 0,45^0 \times 0,55^n = 0,55^n$.

Donc $p_n = 1 - P(X = 0) =  1 - 0,55^n$.
		\item












En partant de $n=0$, le programme calcule $p_n$ et augmente la taille de l'échantillon de 1 tant que $p_n < 0,99$.

		\item 
		On cherche donc $n$ tel que $1 - 0,55^n \geqslant 0,99 \iff 0,01 \geqslant 0,55^n$, d'où par croissance de la fonction logarithme népérien : $\ln 0,01 \geqslant n \ln 0,55 \iff \dfrac{\ln 0,01}{\ln 0,55} \leqslant n$ (car $\ln 001 < 0$).
		
		Or $\dfrac{\ln 0,01}{\ln 0,55} \approx 7,7$.
		
Conclusion : le programme renvoie la valeur 8.
	\end{enumerate}
\end{enumerate}

\bigskip

\end{document}
